\documentclass[11pt, a4paper]{article}

\usepackage[T1]{fontenc}
\usepackage[utf8]{inputenc}
\usepackage{tgpagella}
\usepackage[spanish, es-noshorthands]{babel}
\usepackage{amsmath, amssymb, bm}
\usepackage{booktabs}
\usepackage{tabularx}
\usepackage[table, xcdraw]{xcolor}
\usepackage[most]{tcolorbox}
\usepackage[american]{circuitikz}
\usepackage[margin=2.5cm]{geometry}
\usepackage{fancyhdr}

\pagestyle{fancy}
\fancyhf{}
\setlength{\headheight}{16pt}
\lhead{\textbf{UCB Sede Tarija} | Ing. Mecatrónica}
\chead{}
\rhead{IMT-121 Circuitos Electrónicos I | PATH A}
\lfoot{Prof: M.Sc. Ing. Bernardo Quiroga Turdera}
\rfoot{Página \thepage}
\renewcommand{\headrulewidth}{0.4pt}

\definecolor{ucbblue}{RGB}{0, 51, 102}

\begin{document}

\begin{tcolorbox}[colback=ucbblue!5!white, colframe=ucbblue, title=\textbf{Práctica de Consolidación: Superposición (Path A)}]
    \centering \textbf{Trabajo en Parejas - Superación de Errores Procedimentales}
\end{tcolorbox}

\section*{Parte 1: Clasificación de Errores Comunes}
Lea las siguientes acciones realizadas por un estudiante ficticio al aplicar superposición. Indique qué regla infringió y cuál es la corrección física o matemática.

\begin{enumerate}
    \item \textit{"Al analizar el subcircuito de la fuente de corriente, reemplacé la fuente de voltaje de 10 V por un circuito abierto."} \\ \hrulefill \\ \hrulefill
    
    \item \textit{"Al analizar el subcircuito de la fuente de voltaje, reemplace la fuente de corriente de 5 mA por un cable."} \\ \hrulefill \\ \hrulefill
    
    \item \textit{"La primera contribución fue de $6\,\text{V}$ y la segunda de $-1.2\,\text{V}$. El resultado total que reporté fue $7.2\,\text{V}$ sumando ambas cantidades."} \\ \hrulefill \\ \hrulefill
\end{enumerate}

\section*{Parte 2: Práctica de Cálculo y Signos}

\textbf{Ejercicio A1: Fuentes Aditivas}
En un circuito se tiene:
\begin{itemize}
    \item $V_s = 15\,\text{V}$ (polo positivo conectado a $R_1$).
    \item $R_1 = 5\,\text{k}\Omega$ (entre $V_s$ y el nodo $V_o$).
    \item $R_2 = 10\,\text{k}\Omega$ (entre $V_o$ y tierra).
    \item $I_s = 0.5\,\text{mA}$ (conectada entre tierra y $V_o$, inyectando corriente \textbf{hacia} el nodo $V_o$).
\end{itemize}
Calcule $V_o$ aplicando superposición y prestando extrema atención a las referencias. Muestre sus subcircuitos equivalentes.

\vspace{5cm}

\textbf{Ejercicio A2: Fuentes de Efecto Opuesto}
En un circuito diferente se tiene:
\begin{itemize}
    \item $V_s = 8\,\text{V}$ (polo positivo conectado a $R_1$).
    \item $R_1 = 2\,\text{k}\Omega$ (entre $V_s$ y el nodo $V_o$).
    \item $R_2 = 2\,\text{k}\Omega$ (entre $V_o$ y tierra).
    \item $I_s = 1\,\text{mA}$ (conectada entre el nodo $V_o$ y tierra, extrayendo corriente \textbf{desde} $V_o$ hacia tierra).
\end{itemize}
Calcule $V_o$ aplicando superposición. Compruebe la consistencia de sus signos.

\end{document}
